\chapter{2000年上海卷（春）}
\section{填空题}
\begin{problem}
若 \(\sin \left(\frac{\pi}{2} + \alpha\right) = \frac{3}{5}\)，则 \(\cos 2\alpha =\)\fillinblank{}.
\end{problem}

\begin{answer}
\(-\dfrac{7}{25}\)
\end{answer}

\begin{problem}
若函数 \( f(x) = \frac{x}{x + 2} \)，则 \( f^{-1}\left(\frac{1}{3}\right) = \)\fillinblank{}.
\end{problem}

\begin{answer}
\(1\)
\end{answer}

\begin{problem}
方程 \( \log_{4}(3x-1)=\log_{4}(x-1)+\log_{4}(3+x) \) 的解是\fillinblank{}.
\end{problem}

\begin{answer}
\(x=2\)
\end{answer}

\begin{problem}
若 \( (\sqrt{x}+a)^{5} \) 的展开式中的第四项是 \( 10a^{2} \)（\(a\) 为大于零的常数），则 \(x=\)\fillinblank{}.
\end{problem}

\begin{answer}
\(\dfrac{1}{a}\)
\end{answer}

\begin{problem}
在三角形 \(ABC\) 中，\(\sqrt{2} \sin A = \sqrt{3} \cos A\)，则 \(\angle A =\)\fillinblank{}.
\end{problem}

\begin{answer}
\(60^\circ\)
\end{answer}

\begin{problem}
若直线 \(l\) 的倾角为 \( \pi - \arctan\frac{1}{2} \)，且过点 \( (1,0) \)，则直线 \(l\) 的方程为\fillinblank{}.
\end{problem}

\begin{answer}
\(x+2y-1=0\)
\end{answer}

\begin{problem}
若数列 \(\{a_{n}\}\) 的通项为 \(\frac{1}{n(n + 1)}\)（\(n \in \N\)），则 \(\displaystyle\lim_{n \to \infty} (a_{1} + n^{2} a_{n}) =\)\fillinblank{}.
\end{problem}

\begin{answer}
\(\dfrac{3}{2}\)
\end{answer}

\begin{problem}
如图，\(\angle BAD = 90^{\circ}\) 的等腰直角三角形 \(ABD\) 与正三角形 \(CBD\) 所在平面互相垂直，\(E\) 是 \(BC\) 的
中点，则 \(AE\) 与平面 \(BCD\) 所成角的大小为\fillinblank{}.

\bitmapfigure[max width=0.7\linewidth]{img/2000/shanghai_spring/q8_fig1.png}
\end{problem}

\begin{answer}
\(45^\circ\)
\end{answer}

\begin{problem}
若两个长方体的长，宽，高分别为 \(5\mathrm{~cm}\)，\(4\mathrm{~cm}\)，\(3\mathrm{~cm}\)，把它们两个全等的面重合在一起组成大长方体，则大长方体的对角线最长为\fillinblank{}\(\mathrm{cm}\).
\end{problem}

\begin{answer}
\(5\sqrt{5}\)
\end{answer}

\begin{problem}
有 \(n\)（\(n \in \N\)）件不同的产品排成一排，若其中 \(A\)，\(B\) 两件产品排在一起的不同排法有 48 种，则 \(n = \)\fillinblank{}.
\end{problem}

\begin{answer}
\(5\)
\end{answer}

\begin{problem}
集合 \(A = \{(x, y) | x^2 + y^2 = 4\}\)，\(B = \{(x, y) | (x - 3)^2 + (y - 4)^2 = r^2\}\)，其中 \(r > 0\)，若 \(A \cap B\) 中有且仅有一个元素，则 \(r\) 的值是\fillinblank{}.
\end{problem}

\begin{answer}
\(3\) 或 \(7\)
\end{answer}

\begin{problem}
设 \(I\) 是全集，非空集合 \(P\)，\(Q\) 满足 \(P \subset Q \subset I\). 若含 \(P\)，\(Q\) 的一个集合运算表达式，使运算结果为空集 \(\varnothing\)，则这个运算表达式可以是（只要写出一个表达式）\fillinblank{}.
\end{problem}

\begin{answer}
\(\overline{Q}\cap P\)
\end{answer}

\section{单选题}
\begin{problem}
抛物线 \( y = -x^{2} \) 的焦点坐标为（\quad）

\choices
    {\(\left(0, \frac{1}{4}\right)\).}
    {\(\left(0, -\frac{1}{4}\right)\).}
    {\(\left(\frac{1}{4}, 0\right)\).}
    {\( \left(-\frac{1}{4},0\right) \).}
\end{problem}

\begin{answer}
B
\end{answer}

\begin{problem}
\(x = \sqrt{1 - 3y^2}\) 所表示的曲线是（\quad）

\choices
    {双曲线.}
    {椭圆.}
    {双曲线的一部分.}
    {椭圆的一部分.}
\end{problem}

\begin{answer}
D
\end{answer}

\begin{problem}
“\(a=1\)”是“函数 \( y=\cos^{2}ax-\sin^{2}ax \) 的最小正周期为 \( \pi \) ”的（\quad）

\choices
    {充分不必要条件.}
    {必要不充分条件.}
    {充要条件.}
    {既非充分条件也非必要条件.}
\end{problem}

\begin{answer}
A
\end{answer}

\begin{problem}
若 \(0 < a < 1\)，\(b < -1\)，则函数 \(f(x) = a^x + b\) 的图像不经过（\quad）

\choices
    {第一象限.}
    {第二象限.}
    {第三象限.}
    {第四象限.}
\end{problem}

\begin{answer}
A
\end{answer}
\section{解答题}
\begin{problem}
设 \(f(x)\) 为定义在 \(\R\) 上的偶函数，当 \(x \leqslant -1\) 时，\(y = f(x)\) 的图像是经过点 \((-2,0)\)，斜率为 1 的射线. 又在 \(y = f(x)\) 的图像中有一部分是顶点在 \((0,2)\)，且过点 \((-1,1)\) 的一段抛物线. 试写出函数 \(f(x)\) 的表达式，并作出其图像.

\bitmapfigure[max width=0.48\linewidth]{img/2000/shanghai_spring/q17_fig1.png}
\end{problem}

\begin{solution}
当 \(x\leqslant -1\) 时，设 \(f(x)=x+b\)，则 \(0=-2+b\)，即 \(b=2\)，得 \(f(x)=x+2\)；

当 \(-1<x<1\) 时，设 \(f(x)=ax^2+2\)，则 \(1=a(-1)^2+2\)，即 \(a=-1\)，得 \(f(x)=-x^2+2\)；

当 \(x\geqslant 1\) 时，\(f(x)=-x+2\).

故 \(f(x)=\begin{cases}
x+2,&x\leqslant -1,\\
2-x^2,&-1<x<1,\\
-x+2,&x\geqslant 1.
\end{cases}\)

\solutionfigure{\bitmapfigure[max width=0.42\linewidth]{img/2000/shanghai_spring/q17_solution_fig1.png}}
\end{solution}

\begin{problem}
设复数 \(z\) 满足 \(|z| = 5\)，且 \((3 + 4\i)z\) 在复平面上对应的点在第二，四象限的角平分线上，\(|\sqrt{2}z - m| = 5\sqrt{2}\)（\(m \in \R\)），求 \(z\) 和 \(m\) 的值.
\end{problem}

\begin{solution}
【解法一】设 \(z=x+y\i\)（\(x,y\in\R\)），由 \(|z|=5\)，得 \(x^2+y^2=25\).

又 \((3+4\i)z=(3+4\i)(x+y\i)=(3x-4y)+(4x+3y)\i\)，且 \((3+4\i)z\) 在复平面上对应的点在第二，四象限的角平分线上，所以 \(3x-4y+4x+3y=0\)，得 \(y=7x\).

所以 \(x=\pm\dfrac{\sqrt{2}}{2}\)，\(y=\pm\dfrac{7\sqrt{2}}{2}\)，即 \(z=\pm\left(\dfrac{\sqrt{2}}{2}+\dfrac{7\sqrt{2}}{2}\i\right)\)，从而 \(\sqrt{2}z=\pm(1+7\i)\).

当 \(\sqrt{2}z=1+7\i\) 时，有 \(|1+7\i-m|=5\sqrt{2}\)，即 \((1-m)^2+7^2=50\)，得 \(m=0\) 或 \(m=2\)；当 \(\sqrt{2}z=-(1+7\i)\) 时，同理得 \(m=0\) 或 \(m=-2\).

【解法二】因为 \(|(3+4\i)z|=|3+4\i||z|=25\)，所以 \((3+4\i)z=25\left(\cos\dfrac{3\pi}{4}+\i\sin\dfrac{3\pi}{4}\right)\)，或 \((3+4\i)z=25\left(\cos\dfrac{7\pi}{4}+\i\sin\dfrac{7\pi}{4}\right)\)，得 \(z=\dfrac{\sqrt{2}}{2}+\dfrac{7\sqrt{2}}{2}\i\)，或 \(z=-\left(\dfrac{\sqrt{2}}{2}+\dfrac{7\sqrt{2}}{2}\i\right)\).

当 \(\sqrt{2}z=1+7\i\) 时，有 \((1-m)^2+7^2=50\)，得 \(m=0\) 或 \(m=2\)；当 \(\sqrt{2}z=-(1+7\i)\) 时，同理得 \(m=0\) 或 \(m=-2\).
\end{solution}

\begin{problem}
有一批影碟机（VCD）原销售价为每台 \(800\text{ 元}\)，在甲乙两家家电商场均有销售. 甲商场用如下的方法促销：买一台单价为 \(780\text{ 元}\)，买两台每台单价都为 \(760\text{ 元}\)，依次类推，每多买一台则所买各台单价均再减少 \(20\text{ 元}\)，但每台最低价不能低于 \(440\text{ 元}\)；乙商场一律都按原价的 \(75\%\) 销售. 某单位需购买一批此类影碟机，问去哪家商场购买花费较少？
\end{problem}

\begin{solution}
设某单位需购买 \(x\) 台影碟机，甲，乙两商场的购货款的差为 \(y\).

去甲商场购买共花费 \((800-20x)x\)，根据题意，\(800-20x\geqslant440\)，所以 \(1\leqslant x\leqslant18\)；去乙商场购买共花费 \(600x\)，其中 \(x\in\N\).

所以 \(y=\begin{cases}
(800-20x)x-600x,&1\leqslant x\leqslant18,\\
440x-600x,&18<x
\end{cases}\quad(x\in\N)\)，即 \(y=\begin{cases}
200x-20x^2,&1\leqslant x\leqslant18,\\
-160x,&18<x
\end{cases}\quad(x\in\N)\).

于是 \(\begin{cases}
y>0,&1\leqslant x<10,\\
y=0,&x=10,\\
y<0,&x>10
\end{cases}\quad(x\in\N)\).

故若买少于 \(10\) 台，去乙商场花费较少；若买 \(10\) 台，去甲，乙两商场花费相同；若买超过 \(10\) 台，去甲商场花费较少.
\end{solution}

\begin{problem}
已知 \(\{a_{n}\}\) 是等差数列，\(a_{1} = -393\)，\(a_{2} + a_{3} = -768\)，\(\{b_{n}\}\) 是公比为 \(q\)（\(0 < q < 1\)）的无穷等比数列，\(b_{1} = 2\)，且 \(\{b_{n}\}\) 的各项和为 20.

\begin{enumerate}
\item 写出 \(\{a_{n}\}\) 和 \(\{b_{n}\}\) 的通项公式；
\item 试求满足下列不等式的正整数 \(m\)： \(\displaystyle \frac{a_{m + 1} + a_{m + 2} + \cdots + a_{2m}}{m + 1} \leqslant -160b_{2}\).
\end{enumerate}
\end{problem}

\begin{solution}
\begin{enumerate}
\item 设 \(\{a_n\}\) 的公差为 \(d\)，则 \(a_2+a_3=2a_1+3d\)，所以 \(2\times(-393)+3d=-768\)，解得 \(d=6\). 因而 \(a_n=-393+6(n-1)=6n-399\).

由 \(\dfrac{b_1}{1-q}=20\)，得 \(\dfrac{2}{1-q}=20\)，所以 \(q=\dfrac9{10}\)，从而 \(b_n=2\left(\dfrac9{10}\right)^{n-1}\)（\(n\in\N\)）.

\item 因为 \(a_1+a_2+\cdots+a_m=ma_1+\dfrac{m(m-1)d}{2}=-393m+3m(m-1)\)，所以
\(a_{m+1}+a_{m+2}+\cdots+a_{2m}=(a_1+a_2+\cdots+a_{2m})-(a_1+a_2+\cdots+a_m)=9m^2-396m\).

又 \(-160b_2=-288\)，所以 \(9m^2-396m\leqslant-288(m+1)\)，即 \(m^2-44m\leqslant-32(m+1)\)，得 \((m-4)(m-8)\leqslant0\). 因为 \(m\in\N\)，所以 \(m=4,5,6,7,8\).
\end{enumerate}
\end{solution}

\begin{problem}
四棱锥 \(P\)—\(ABCD\) 中，底面 \(ABCD\) 是一个平行四边形，\(\overrightarrow{AB}=\{2,-1,-4\}\)，\(\overrightarrow{AD}=\{4,2,0\}\)，\(\overrightarrow{AP}=\{-1,2,-1\}\).

\begin{enumerate}
\item 求证：\(PA \perp\) 底面 \(ABCD\)；
\item 求四棱锥 \(P\)—\(ABCD\) 的体积；
\item 对于向量 \(\bs{a}=\{x_{1},y_{1},z_{1}\}\)，\(\bs{b}=\{x_{2},y_{2},z_{2}\}\)，\(\bs{c}=\{x_{3},y_{3},z_{3}\}\)，定义一种运算：\(\displaystyle (\bs{a} \times \bs{b}) \cdot \bs{c} = x_{1} y_{2} z_{3} + x_{2} y_{3} z_{1} + x_{3} y_{1} z_{2} - x_{1} y_{3} z_{2} - x_{2} y_{1} z_{3} - x_{3} y_{2} z_{1}\). 试计算 \( (\overrightarrow{AB}\times\overrightarrow{AD})\cdot\overrightarrow{AP} \) 的绝对值的值；说明其与四棱锥 \(P\)—\(ABCD\) 体积的关系，并由此猜想向量这一运算 \( (\overrightarrow{AB}\times\overrightarrow{AD})\cdot\overrightarrow{AP} \) 的绝对值的几何意义.
\end{enumerate}
\end{problem}

\begin{solution}
\begin{enumerate}
\item \(\overrightarrow{AP}\cdot\overrightarrow{AB}=-2-2+4=0\)，所以 \(AP\perp AB\). 又 \(\overrightarrow{AP}\cdot\overrightarrow{AD}=-4+4+0=0\)，所以 \(AP\perp AD\). 因为 \(AB\)，\(AD\) 是底面 \(ABCD\) 上的两条相交直线，所以 \(AP\perp\) 底面 \(ABCD\).

\item 设 \(\overrightarrow{AB}\) 与 \(\overrightarrow{AD}\) 的夹角为 \(\theta\)，则 \(\displaystyle \cos\theta=\frac{\overrightarrow{AB}\cdot\overrightarrow{AD}}{|\overrightarrow{AB}||\overrightarrow{AD}|}=\frac{8-2}{\sqrt{4+1+16}\sqrt{16+4}}=\frac3{\sqrt{105}}\).

所以四棱锥体积 \(\displaystyle V=\frac13|\overrightarrow{AB}||\overrightarrow{AD}|\sin\theta|\overrightarrow{AP}|=\frac23\sqrt{105}\sqrt{1-\frac9{105}}\sqrt{1+4+1}=16\).

\item \(\displaystyle |(\overrightarrow{AB}\times\overrightarrow{AD})\cdot\overrightarrow{AP}|=|-4-32-4-8|=48\)，它是四棱锥 \(P\text{--}ABCD\) 体积的 \(3\) 倍.

由此猜想，\(|(\overrightarrow{AB}\times\overrightarrow{AD})\cdot\overrightarrow{AP}|\) 的几何意义是以 \(\overrightarrow{AB}\)，\(\overrightarrow{AD}\)，\(\overrightarrow{AP}\) 为棱的平行六面体的体积（或以它们为棱的直四棱柱的体积）.
\end{enumerate}
\end{solution}

\begin{problem}
如图所示， \(A\)，\(F\) 分别是椭圆 \(\frac{(y + 1)^2}{16} +\frac{(x - 1)^2}{12} = 1\) 的一个顶点与一个焦点，位于 \(x\) 轴的正半轴上的动点 \(T(t,0)\) 与 \(F\) 的连线交
射线 \(OA\) 于 \(Q\). 求：

\begin{enumerate}
\item 点 \(A\)，\(F\) 的坐标及直线 \(TQ\) 的方程；
\item \(\triangle OTQ\) 的面积 \(S\) 与 \(t\) 的函数关系式 \(S = f(t)\) 及该函数的最小值；
\item 写出 \(S = f(t)\) 的单调递增区间，并证明之.
\end{enumerate}

\bitmapfigure[max width=0.48\linewidth]{img/2000/shanghai_spring/q22_fig1.png}
\end{problem}

\begin{solution}
\begin{enumerate}
\item 点 \(A\) 的坐标为 \((1,3)\)，点 \(F\) 的坐标为 \((1,1)\).

当 \(t>0\) 且 \(t\neq1\) 时，直线 \(TQ\) 的方程为 \(\displaystyle y=\frac{x-t}{1-t}\)；当 \(t=1\) 时，直线 \(TQ\) 的方程为 \(x=1\).

\item 联立直线 \(OA\) 和直线 \(TQ\) 的方程：当 \(t\neq1\) 时，\(\displaystyle \begin{cases}
y=3x,\\
y=\dfrac{x-t}{1-t}
\end{cases}\)，当 \(t=1\) 时，\(\displaystyle \begin{cases}
y=3x,\\
x=1
\end{cases}\). 得 \(\displaystyle y_Q=\frac{3t}{3t-2}\)（\(t\neq1\)），或 \(y_Q=3\)（\(t=1\)）.

因为 \(Q\) 在射线 \(OA\) 上，所以 \(t>0\) 且 \(y_Q>1\)，从而 \(t>\dfrac23\).

于是 \(\displaystyle S=f(t)=\frac12ty_Q=\frac{3t^2}{2(3t-2)}\)（\(t>\dfrac23\)）.

又 \(\displaystyle f(t)=\frac{3t^2}{2(3t-2)}=\frac{(3t)^2}{6(3t-2)}=\frac{(3t-2+2)^2}{6(3t-2)}=\frac16\left(3t-2+\frac4{3t-2}+4\right)\). 因为 \(3t-2>0\)，所以 \(\displaystyle f(t)\geqslant\frac16(2\sqrt4+4)=\frac43\)，当且仅当 \(3t-2=2\)，即 \(t=\dfrac43\) 时等号成立. 因此 \(S=f(t)\) 的最小值为 \(\dfrac43\).

\item \(S=f(t)=\dfrac{3t^2}{2(3t-2)}\) 在区间 \(\left(\dfrac43,+\infty\right)\) 上为增函数.

证明：任取 \(t_1,t_2\in\left(\dfrac43,+\infty\right)\)，不妨设 \(t_2>t_1>\dfrac43\)，则
\[
\begin{aligned}
f(t_1)-f(t_2)
&=\frac16\left(3t_1-2+\frac4{3t_1-2}+4\right)-\frac16\left(3t_2-2+\frac4{3t_2-2}+4\right) \\
&=\frac12(t_1-t_2)\left[1-\frac4{(3t_1-2)(3t_2-2)}\right] \\
&=\frac12(t_1-t_2)\frac{(3t_1-2)(3t_2-2)-4}{(3t_1-2)(3t_2-2)}.
\end{aligned}
\]

因为 \(t_2>t_1>\dfrac43\)，所以 \(t_1-t_2<0\)，\((3t_1-2)(3t_2-2)>4\)，从而 \(f(t_1)<f(t_2)\)，故 \(S=f(t)\) 在 \(\left(\dfrac43,+\infty\right)\) 上为增函数.
\end{enumerate}
\end{solution}
